\section{Implementation and independently replayed certificates}
\label{sec:implementation}

The delivered extension is opt-in and uses the Python standard library.
The pinned native runtime is included for reproduction.  Regina is used
only by the explicit-cut oracle, and is not imported by the new producer
or verifier.  The default recognizer and its schedules are unchanged.

\subsection{The producer}

\texttt{normal\_cut\_saturation} accepts the same source triangulation and
normal-coordinate contract as the existing normal-cut parallelity operator.
Its logical sequence is:
\begin{enumerate}
\item Construct the source parallelity decomposition, with four orbit
queries and optional upstream certificates.  Reject unsupported core
three-handles or closed bundle bases as inconclusive.
\item If the caller enables the original-core precheck, form its fixed
forest and capped source incidence rows.  A nonempty stopping core permits
an early inconclusive result; no ordered boundary information is needed.
\item Obtain the ordered attachment skeleton, requiring at most two more
orbit queries.  Source-zero-port cases may require none.
\item Compile all original core rows and all saturated residual rows
jointly against the original source cocores.  Record the exact private
one-handle count in each connected attachment component.
\item Run unit-column peeling.  Complete elimination authorizes the
handlebody histogram and Euler check.  A nonempty residual authorizes only
an inconclusive result for topology.
\end{enumerate}

The discovery cycle allowance is shared across the upstream queries.
The optional precheck defaults to false, so the principal audit and
certificate format are directly comparable to the complete six-query
path.  Its successful path returns the same whole-cut certificate as the
ordinary path; its early obstruction uses a smaller four-query schema.

\subsection{What is and is not trusted}

The whole-cut certificate has four proof layers and an exact conclusion:
\begin{center}\small
\begin{tabularx}{\linewidth}{@{}p{0.21\linewidth}YY@{}}
\toprule
Layer & Claimed data & Independent check\\
\midrule
Source parallelity & Core handles, bundle bases, fibre cover and port
ownership & Reconstruct the source geometry and verify the four orbit proofs\\
Ordered boundary & Maximal interval/circle interfaces and signed source
port walks & Independently reconstruct and replay the endpoint and gap data\\
Relative compiler & Forest, residual rows, attachments, component mergers,
private counts & Rebuild finite source paths and relative counts without
calling the compiler\\
Peeling & Row/column deletion schedule and final residual & Timestamp scan
of immutable source incidences; stopping inequalities\\
Conclusion & Exact genus histogram or exact raw-obstruction statement &
Recompute ranks, multiplicities, Euler characteristic, and result fields\\
\bottomrule
\end{tabularx}
\end{center}

The enclosing verifier imports neither the relative compiler nor the
peeling producer.  The existing source verifiers reconstruct the source
geometry independently rather than accepting a geometry dictionary because
it carries a matching checksum.  Both source proofs are bound to the
\emph{same} input triangulation and normal vector.  The finite compiler
verifier alone is conditional on its supplied geometric inputs; it is the
enclosing verifier that closes this source-authentication boundary.

A source fingerprint guards against accidental substitution, but it is not
the mathematical reason a conclusion is correct.  The reason is the
reconstruction, the geometric theorems, and the exact incidence replay.
The certificate includes no unchecked Boolean saying that a residual row
is geometrically valid or that a component is a ball.

The early certificate follows the same principle with fewer layers.  It
contains the authenticated source parallelity proof, a finite core forest
and row record, a peeling certificate with nonempty residual, and the exact
conclusion \texttt{INCONCLUSIVE}.  Its checker independently reconstructs
the core rows before verifying the stopping set.  By
Corollary~\ref{cor:row-subset-monotonicity}, no eventual saturated row can
make that fixed core subset peel.

\subsection{Status and integer semantics}

\texttt{READY} in the finite compiler means that a relative presentation
has been compiled.  It is not a topology result.  \texttt{COMPLETE} in the
whole-cut consumer means that every row has cancelled and that the exact
handlebody census follows.  \texttt{OBSTRUCTED} in the matrix layer means
that full-row raw positive cancellation is impossible.  At the topology
boundary it becomes \texttt{INCONCLUSIVE}.

The source matrices retain empty rows.  Omitting an empty residual row
would incorrectly turn the M\"obius-on-a-ball example into a ball
certificate.  Counts are capped individually at two only for incidence
logic.  The total of those capped counts remains uncapped for deletion
bookkeeping; saturating the total itself would lose the information needed
to know when a column becomes a leaf.

Encoded integers use the existing native hexadecimal convention for very
large values.  The code never disables Python's global decimal conversion
limit.  Genus, multiplicity, and Euler arithmetic use exact integers.
Boolean values are rejected where exact integers are required.  Certificate
record keys and canonical lists are checked, and caller cancellation
exceptions propagate rather than being swallowed as false proofs.

\subsection{Why integration remains opt-in}

Replacing the old minimum-fill routine globally would change the meaning
of its partial output.  On an obstructed matrix, the old routine can make
many legal partial cancellations, whereas column peeling may stop
immediately with a canonical residual.  Neither method promises maximum
partial cancellation.  A consumer using only the full-row verdict can
switch safely after the stated theorem and tests; a consumer using the
old partial schedule needs a separate interface review.

The supplied integration patch therefore adds new modules without changing
the default solver or its old proof schemas.  A caller can explicitly
invoke the extended whole-cut consumer and independently replay its result.
Promoting it into a knot-recognition schedule would additionally require
the selected-surface and peripheral correctness obligations of that
schedule.  The integration notes list the exact files and test commands.
